\begin{activity}{Half-Space Cooling of the Oceanic Lithosphere}

\section*{Purpose}
To develop physical intuition for transient cooling of the oceanic lithosphere
using the half-space cooling model.

\section*{Model Assumptions}
\begin{itemize}
  \item The mantle is initially at uniform temperature.
  \item Surface temperature is suddenly reduced and held fixed.
  \item No fixed lithospheric thickness is imposed.
\end{itemize}

\section*{Solution Form}
The temperature evolution is given by
\[
T(z,t) = T_m \,\mathrm{erf}\!\left(
\frac{z}{2\sqrt{\kappa t}}
\right).
\]

\section*{Tasks}
\begin{enumerate}
  \item Sketch temperature--depth profiles for increasing times.
  \item Identify the thermal boundary layer and describe how its thickness evolves.
  \item Estimate how the depth to a given isotherm scales with seafloor age.
\end{enumerate}

\section*{Geological Interpretation}
\begin{itemize}
  \item Oceanic lithosphere thickens thermally with age.
  \item Heat flow decreases as the plate cools away from mid-ocean ridges.
\end{itemize}

\section*{Key Insight}
Explain why the half-space cooling model predicts unbounded lithospheric thickening.

\end{activity}